% !TEX root = ../../../main.tex
% Appendix entry: detection.wavelet_summation -- "Wavelet summation" (Appendix A).
\entryfield{Algorithm}
Stage~2 of the \texttt{dehshibi\_spikes} template (stage~1 is
\blockref{preprocessing.wavelet_transform}, stage~3 \blockref{detection.summation_threshold}).
\code{_run} in \code{Adapters/detection_wavelet_summation.py}, for an image $G$ with $R$ rows
(row~0 the lowest frequency) and fractions $f_0 = \texttt{row\_from}$, $f_1 = \texttt{row\_to}$:
\begin{enumerate}
\item $r_0 = \mathrm{round}(f_0 R)$, $r_1 = \mathrm{round}(f_1 R)$ (\code{band_rows});
  refuse if $r_1 \le r_0$.
\item $\Omega(\tau) = \sum_{r = r_0}^{r_1 - 1} G(r, \tau)$, summed in float64.
\end{enumerate}
In the authors' code (\code{iApplyWavelet.m}, method~2) $\Omega$ is the sum over the rows whose
frequency is at most a quarter of the frequency range, not over every scale as the printed paper
reads \citep{dehshibi2021electrical}. For a 3000-sample window that is the lowest 66 of 87 rows,
which is why the template passes $\texttt{row\_to} = 0.7586$. The band is a fraction because an
\texttt{Encoding} carries no frequency axis and a block cannot read the previous block's
\code{meta}; the transform prints the fraction for any other window.

\emph{The contract.} Any block that emits an \texttt{Encoding} of kind \texttt{image} with one
column per sample can feed this one; any block that emits a per-sample \texttt{Scores} can
replace it in front of \blockref{detection.summation_threshold}
(\code{docs/BLOCK_INTEGRATION.md}, \S6).

\entryfield{Assumptions}
\begin{itemize}
\item \emph{Units.} None: the values are the transform's dimensionless counts (1--241 per row).
\item \emph{Shape.} The image must have exactly one column per sample of the span; anything
  else is refused. \blockref{detection.wavelet_scattering} and \code{detection.freq_stft} emit
  one column per frame, not per sample, and cannot feed it.
\item \emph{Missing values.} A NaN anywhere in a column makes that column's $\Omega$ NaN; the
  count is in \code{meta["n_nan"]}. The next block skips any window whose score is not finite
  throughout.
\end{itemize}

\entryfield{Outputs}
A \texttt{Scores} with one float per sample and the span's $f_s$. \code{meta} holds
\code{n_rows}, \code{rows_used} ($[r_0, r_1)$) and \code{n_nan}.

\entryfield{Cost}
A column sum: $O(Rn)$ time and $O(n)$ extra memory beyond the input image. It declares no
estimate and no span cap; the cap on the image upstream bounds it.

\entryfield{Failure modes}
\begin{itemize}
\item \emph{The wrong band, silently.} The template's 0.7586 is the authors' band only when
  the transform made 87 rows (a 3000-sample window). At 10\,Hz, or with another
  \emph{window\_s}, the image has a different number of rows and the fixed fraction sums a
  different band from the authors'. Nothing warns: read the transform's readout and copy its
  fraction.
\item \emph{Alone, it sums everything.} Its own defaults (0 to 1) reproduce the printed paper's
  reading, not the authors' code.
\item \emph{Rounding.} The band edges are rounded to whole rows, so two fractions closer than
  $1/R$ can give the same band.
\end{itemize}

\entryfield{Validation}
\code{tests/test_adapter_wavelet_summation.py} checks the sum, the band, NaN propagation and
the refusals, and that the template's band applied to the transform's image gives the
$\Omega$ the authors' MATLAB computed on the fixture recording, to within one count
(\code{test_the_templates_band_gives_the_authors_omega}). That check is at 1\,Hz on a
3000-sample recording only.

\entryfield{References}
\citet{dehshibi2021electrical}; the authors' code, Zenodo doi:10.5281/zenodo.3997031;
\code{docs/prompts/fixup/reports/J-dehshibi-vs-the-paper.md}, \S4 (row ``$\Omega$'').

\entryfield{Tests}
\code{tests/test_adapter_wavelet_summation.py}; \code{tests/test_template_dehshibi.py} runs it
inside the template.
